% !TEX root = ../../../main.tex
\section{Cauchy criterion for uniform convergence}
\label{sec:analysis-i:cauchy-criterion-for-uniform-convergence}

% Editorial completion authorised by the HW06 topic-coverage request.
\begin{pendingreview}
The Cauchy criterion tests convergence without first identifying the
limit. For functions, the comparisons must hold at every input with
the same index threshold.

\begin{definition}[Uniformly Cauchy sequence]
  \label{thm:analysis-i:uniformly-cauchy}
  A sequence $f_n:D\to\R$ is \emph{uniformly Cauchy} if
  \[
    \begin{gathered}
      \forall\varepsilon>0\quad\exists N\in\N\\
      \forall m,n\ge N\quad\forall x\in D\quad
      |f_m(x)-f_n(x)|<\varepsilon.
    \end{gathered}
  \]
\end{definition}

\begin{theorem}[Uniform Cauchy criterion]
  \label{thm:analysis-i:uniform-cauchy-criterion}
  A sequence of real-valued functions on $D$ converges uniformly if
  and only if it is uniformly Cauchy.
\end{theorem}
\begin{proof}
  If $f_n\to f$ uniformly, choose $N$ so that
  $|f_n(x)-f(x)|<\varepsilon/2$ for every $n\ge N$ and $x\in D$.
  Then the triangle inequality gives
  $|f_m(x)-f_n(x)|<\varepsilon$ for all $m,n\ge N$ and $x\in D$.

  Conversely, suppose $(f_n)$ is uniformly Cauchy. At each fixed $x$,
  the sequence $(f_n(x))$ is Cauchy in $\R$, so completeness gives a
  finite limit $f(x)$. Given $\varepsilon>0$, choose $N$ from the
  uniform Cauchy property with tolerance $\varepsilon/2$. For fixed
  $n\ge N$ and $x\in D$, let $m\to\infty$ in that estimate to get
  \[
    |f(x)-f_n(x)|\le\frac{\varepsilon}{2}<\varepsilon.
  \]
  The index $N$ is independent of $x$, which proves uniform convergence.
\end{proof}

For a sequence in $C[a,b]$, the condition is precisely that
$\|f_m-f_n\|_\infty\to0$ as $m,n\to\infty$. The limit belongs to
$C[a,b]$ by \autoref{thm:analysis-i:uniform-limit-continuity}.
Thus Cauchy sequences for the supremum norm have limits in this same
normed space.

\subsection{The criterion for a series}
Apply \autoref{thm:analysis-i:uniform-cauchy-criterion} to the partial
sums in \autoref{thm:analysis-i:function-series}. The series
$\sum_{n=0}^{\infty}f_n$ converges uniformly on $D$ exactly when
\[
  \begin{gathered}
    \forall\varepsilon>0\quad\exists N\in\N\\
    \forall m>n\ge N\quad\forall x\in D\quad
    \left|\sum_{k=n+1}^{m}f_k(x)\right|<\varepsilon.
  \end{gathered}
\]
The small quantity is the sum of a whole tail, rather than just one
term of the series.

\subsection{The Weierstrass majorant test}
\begin{theorem}[Weierstrass $M$-test]
  \label{thm:analysis-i:weierstrass-m-test}
  Suppose $f_n:D\to\R$ and there are numbers $M_n\ge0$ such that
  \[
    \forall x\in D\quad |f_n(x)|\le M_n,\qquad
    \sum_{n=0}^{\infty}M_n<\infty.
  \]
  Then both $\sum f_n$ and $\sum |f_n|$ converge uniformly on $D$.
  In particular, $\sum f_n(x)$ converges absolutely at every input.
  Writing $S=\sum f_n$ and $S_N=\sum_{n=0}^N f_n$, we have
  \[
    \forall x\in D\quad
    |S(x)-S_N(x)|\le\sum_{n=N+1}^{\infty}M_n.
  \]
\end{theorem}
\begin{proof}
  The convergent numerical series $\sum M_n$ has arbitrarily small
  tails. For $m>n$ and every $x\in D$,
  \[
    \left|\sum_{k=n+1}^{m}f_k(x)\right|
      \le\sum_{k=n+1}^{m}|f_k(x)|
      \le\sum_{k=n+1}^{m}M_k.
  \]
  The same estimate holds with $f_k$ replaced by $|f_k|$. The uniform
  Cauchy criterion therefore proves both uniform-convergence claims.
  Letting $m\to\infty$ in the estimate with $n=N$ gives the stated
  remainder bound. Absolute convergence at each input also follows
  from comparison with $\sum M_n$.
\end{proof}

The majorant $M_n$ must be independent of $x$ on the whole domain under
consideration. Once it is found, a numerical tail estimate controls all
functional errors at once. If every $f_n$ is continuous on $[a,b]$,
its partial sums are continuous. The sum is then continuous by
\autoref{thm:analysis-i:uniform-limit-continuity}.

\begin{example}[A uniform geometric tail]
  For $0<r<1$ and $x\in[-r,r]$, we have $|x^n|\le r^n$. The $M$-test
  proves absolute and uniform convergence of $\sum_{n=0}^{\infty}x^n$
  there, with the explicit estimate
  \[
    \left|\frac1{1-x}-\sum_{n=0}^{N}x^n\right|
      \le\frac{r^{N+1}}{1-r}.
  \]
  On the larger interval $(-1,1)$ the series is pointwise absolutely
  convergent but is not uniformly convergent. Uniform convergence
  would force $x^n$ to tend uniformly to zero, whereas
  $\sup_{|x|<1}|x^n|=1$ for every $n$.
\end{example}
\end{pendingreview}
